% TOP_PROBLEM: {"id": 2307061, "title": "Research Problems in Function Theory — Problem 7.61", "category_id": 8, "category": {"id": 8, "name": "analysis", "display_name": "Analysis", "description": "Limits, continuity, calculus, and function theory.", "slug": "analysis", "order_index": 8, "created_at": "2026-07-31T15:26:25.671Z"}, "status": "open"}
% FIRST_POSTED: null
% ENTRY_KIND: statement_only
% Imported from: batch06.tex; UnsolvedMath ID 2307061
% Compile twice: lualatex 2307061.tex
\documentclass[11pt]{article}
\usepackage{fontspec}
\setmainfont{Latin Modern Roman}
\usepackage{amsmath,amssymb,mathtools,unicode-math}
\setmathfont{Latin Modern Math}
\usepackage{xurl,hyperref}
\hypersetup{hidelinks,pdftitle={UnsolvedMath Problem 2307061}}
\newcommand{\sref}[2]{\hyperlink{source:#1:#2}{[S#2]}}
\newcommand{\cataloguescope}[1]{\paragraph{Mathematical question.}#1}
\begin{document}
% UnsolvedMath ID 2307061; source catalogue code AMR-022-7061
\setcounter{section}{2307060}
\section{Research Problems in Function Theory — Problem 7.61}\label{2307061}
{\small\sffamily UnsolvedMath ID 2307061 · Analysis}
\subsection{Definitions and mathematical statement}

\paragraph{Shared notation.}$\mathbb N=\{1,2,\ldots\}$, and $\mathbb Z,\mathbb Q,\mathbb R,\mathbb C$ denote the integers, rationals, reals and complex numbers. Any other conventions are specified locally.


For $n\ge1$ let $\mathcal E=\mathcal O(\mathbb C^n)$ be all entire functions. For a multi-index $\alpha=(\alpha_1,\ldots,\alpha_n)\in\{0,1,\ldots\}^n$, put $z^\alpha=\prod z_j^{\alpha_j}$ and $D^\alpha=\prod(\partial/\partial z_j)^{\alpha_j}$. If $P(z)=\sum_\alpha p_\alpha z^\alpha$ is a polynomial, define $P(D)=\sum_\alpha p_\alpha D^\alpha$ and $T_{P,Q}f=P(D)(Qf)$. Entire exponential type means $|f(z)|\le A e^{B|z|}$ for some finite positive $A,B$.Define $P^*(z)=\overline{P(\bar z)}=\sum_\alpha\overline{p_\alpha}z^\alpha$. The zero polynomial makes $T_{P^*,P}$ the zero operator; the substantive bijection question is for nonzero $P$, including nonzero constants.
\paragraph{Statement recorded in the supplied source.}
The original asks whether $T_{P^*,P}:\mathcal E\to\mathcal E$ is bijective for every polynomial $P$. For $P=0$ the answer is negative, since the operator is zero. For every nonzero complex polynomial $P$, is this operator bijective? Also consider the source's proposed sufficient route: for every multi-index $\alpha$, does
\[P^*(D)(Pf)=z^\alpha
\]
admit an entire solution of exponential type? The source records the homogeneous case as known and notes the injectivity claim separately; neither replaces the full nonhomogeneous bijection question. \sref{2307061}{2}
\subsection{Short English statement}
Determine bijectivity of polynomial multiplication followed by the differential operator with conjugate coefficients, and whether exponential-type solutions exist for every monomial right-hand side.
\subsection{Sources}
\begin{description}
\item[\hypertarget{source:2307061:1}{S1}] ULAM.AI and the UnsolvedMath Contributors, UnsolvedMath: A Curated Collection of Open Mathematics Problems, record 2307061 (AMR-022-7061). CC BY 4.0. \url{https://huggingface.co/datasets/ulamai/UnsolvedMath}.
\item[\hypertarget{source:2307061:2}{S2}] Walter K. Hayman and Edward F. Lingham, Research Problems in Function Theory (2018), Problem 7.61, its complete update and relevant preceding definitions. \url{https://arxiv.org/abs/1809.07200}.
\end{description}
\label{end:2307061}
\end{document}
